\documentclass[a4paper,11pt]{article}
\usepackage[a4paper,margin=2cm]{geometry}
\usepackage[T1]{fontenc}
\usepackage[utf8]{inputenc}
\usepackage[ngerman,english]{babel}
\usepackage{lmodern}
\usepackage{microtype}
\usepackage[table]{xcolor}
\usepackage{parskip}
\usepackage{enumitem}
\usepackage[most]{tcolorbox}
\usepackage{titlesec}
\usepackage{xparse}
\usepackage[hidelinks]{hyperref}
\definecolor{HeaderBlueDark}{RGB}{70,110,170}
\definecolor{RowBlueLight}{RGB}{235,243,255}
\definecolor{RowBlueDark}{RGB}{220,233,250}
\definecolor{SoftGray}{RGB}{90,90,90}
\setlength{\parindent}{0pt}
\setlist[itemize]{leftmargin=1.4em,itemsep=0.35em,topsep=0.35em}
\linespread{1.06}
\titleformat{\section}[block]{\Large\bfseries\color{HeaderBlueDark}}{}{0pt}{}
\titlespacing{\section}{0pt}{1.3em}{0.8em}
\titleformat{\subsection}[block]{\large\bfseries\color{HeaderBlueDark}}{}{0pt}{}
\titlespacing{\subsection}{0pt}{1.0em}{0.6em}
\newtcolorbox{exbox}{enhanced,breakable,colback=RowBlueLight,colframe=RowBlueDark,boxrule=0.4pt,arc=1.5mm,left=1.2mm,right=1.2mm,top=1mm,bottom=1mm,title={Beispiele \textnormal{(Examples)}},fonttitle=\bfseries,colbacktitle=RowBlueDark,coltitle=black}
\NewDocumentEnvironment{grammarentry}{mm+b}{%
\begin{tcolorbox}[enhanced,breakable,colback=white,colframe=HeaderBlueDark,boxrule=0.6pt,arc=1.5mm,left=1.5mm,right=1.5mm,top=1mm,bottom=1mm,colbacktitle=HeaderBlueDark,coltitle=white,fonttitle=\bfseries,title={#1\hfill\normalfont\footnotesize #2}]
#3
\end{tcolorbox}}{}
\newcommand{\GermanText}[1]{\textbf{Deutsch: }#1\par}
\newcommand{\EnglishText}[1]{{\small\color{SoftGray}\textbf{English: }#1\par}}
\newcommand{\ExampleItem}[2]{\item \textbf{#1}\\{\small\color{SoftGray}#2}}
\title{unter ... hindurch\\ \large unter + dative phrase + hindurch}
\date{}
\begin{document}
\maketitle
\tableofcontents
\newpage
\section{Einordnung und Wortart / Classification and part of speech}
\begin{grammarentry}{unter ... hindurch}{Grammatische Einordnung / grammatical classification}
\GermanText{Grammatischer Status: Zirkumpositionsartige Orts- und Richtungsverbindung.}
\EnglishText{Grammatical status: Circumposition-like place and direction pattern.}
\end{grammarentry}
\section{Struktur, Stellung und Rektion / Structure, position and case}
\begin{grammarentry}{Klammerbau / Framing structure}{Kasus / case}
\GermanText{Muster: unter + Dativgruppe + hindurch. Rektion: Dativ (für die hier gemeinte Lage des Weges).}
\EnglishText{Pattern: unter + dative phrase + hindurch. Case: dative (for the path located underneath).}
\end{grammarentry}
\section{Bedeutung / Meaning}
\begin{grammarentry}{Grundbedeutung / Core meaning}{Semantik / semantics}
\GermanText{beschreibt das Passieren einer Stelle auf einem Weg unter einem Hindernis oder Bauwerk hindurch.}
\EnglishText{describes passing through an area underneath an obstacle or structure.}
\end{grammarentry}
\section{Verwendung und Register / Usage and register}
\begin{grammarentry}{Gebrauch / Usage}{Typische Kontexte / typical contexts}
\GermanText{Hier bezeichnet unter die räumliche Lage des Weges (unter der Brücke). Deshalb Dativ, obwohl die gesamte Situation eine Bewegung beschreibt.}
\EnglishText{Here unter locates the path (under the bridge), so dative is used even though the whole event involves movement.}
\end{grammarentry}
\section{Beispiele / Examples}
\begin{exbox}
\begin{itemize}
\ExampleItem{Der Zug fährt unter der Brücke hindurch.}{The train passes under the bridge.}
\ExampleItem{Wir liefen unter einem Torbogen hindurch.}{We walked through beneath an archway.}
\ExampleItem{Das Wasser strömt unter dem Weg hindurch.}{The water flows beneath the path.}
\ExampleItem{Sie duckte sich und ging unter dem Ast hindurch.}{She ducked and went under the branch.}
\end{itemize}
\end{exbox}
\section{Abgrenzung / Comparison with related forms}
\begin{grammarentry}{Vergleich / Comparison}{Unterschiede / differences}
\GermanText{Unter die Brücke fahren (Akkusativ) kann eine Bewegung zu einem Ort unter der Brücke bezeichnen; unter der Brücke hindurchfahren beschreibt das Durchqueren darunter.}
\EnglishText{Unter die Brücke fahren (accusative) can mean driving to a position under the bridge; unter der Brücke hindurchfahren means passing all the way beneath it.}
\end{grammarentry}
\section{Typischer Fehler / Typical mistake}
\begin{grammarentry}{Kasus und Form / Case and form}{Falsch versus richtig / wrong versus correct}
\GermanText{Falsch: unter die Brücke hindurchfahren. Richtig: unter der Brücke hindurchfahren.}
\EnglishText{Incorrect: unter die Brücke hindurchfahren. Correct: unter der Brücke hindurchfahren. Use dative to locate the passage beneath the bridge in this pattern.}
\end{grammarentry}
\section{Kurze Übung / Short exercise}
\begin{grammarentry}{Ergänze die richtige Form / Complete the phrase}{Selbstkontrolle / self-check}
\GermanText{Ergänze die Lücke: Wir fuhren unter ... (die Brücke) hindurch.}
\EnglishText{Fill in the blank: Wir fuhren unter ... (die Brücke) hindurch.}
\end{grammarentry}
\subsection{Lösung / Answer}
\begin{grammarentry}{Lösung / Answer}{Kasus prüfen / check the case}
\GermanText{Wir fuhren unter der Brücke hindurch.}
\EnglishText{We drove under the bridge.}
\end{grammarentry}
\section{Zusammenfassung / Summary}
\begin{grammarentry}{Merksatz / Key point}{Form und Bedeutung / form and meaning}
\GermanText{unter ... hindurch — Dativ (für die hier gemeinte Lage des Weges). beschreibt das Passieren einer Stelle auf einem Weg unter einem Hindernis oder Bauwerk hindurch.}
\EnglishText{unter ... hindurch — dative (for the path located underneath). describes passing through an area underneath an obstacle or structure.}
\end{grammarentry}
\section{Quellen / Sources}
\begin{grammarentry}{Fachliche Einordnung / Classification}{IDS und Duden / IDS and Duden}
\GermanText{Für die Abgrenzung zwischen eindeutigen Zirkumpositionen und ähnlichen Konstruktionen siehe IDS (grammis).}
\EnglishText{For the distinction between unambiguous circumpositions and similar structures, see IDS (grammis).}
\url{https://grammis.ids-mannheim.de/terms/view/506}
\url{https://grammis.ids-mannheim.de/systematische-grammatik/1422}
\end{grammarentry}
\end{document}
